\documentclass{article}

\title{Ch. 4.1. - 4.3. HW Solns}
\author{Mr. Spence}
\date{\today}

\begin{document}

\maketitle

\subsection*{QuickQuiz}

\begin{enumerate}

\item (c) -- A market doesn't have to be in a single location like options (a), (b), and (d).

\item (b) -- Sellers are price-takers in perfectly competitive markets.

\item (a) -- Eggs are all mostly the same regardless of the seller, and there are far too many buyers and sellers for any individual to affect the price. 

\item (b) -- A change in the price of hamburgers leads to a movement along the demand curve.

\item (a) -- When the price of a substitute goes up, the demand for pizza increases.

\item (d) -- Demand for an inferior good falls when incomes rise.

\end{enumerate}

\subsection*{Questions for Review}

\begin{enumerate}
\setcounter{enumi}{2}

\item A change in tastes leads to a shift in the curve because a non-price determinant of quantity demanded is changing. That means the quantity demanded at any given price will change after the change in taste, so the curve must change as well. A change in price leads to a movement along the curve because the curve already shows how quantity demanded changes with price. 

\item Pumpkin juice is an inferior good because his demand for it rises as his income falls. His demand curve shifts right because he demands more. 

\end{enumerate}

\end{document}